\begin{abstract}
Media bias detection typically relies on content analysis or costly manual annotation.
We investigate whether user interaction patterns on social news platforms can serve as complementary signals for perceived media bias.
We construct a large-scale dataset from Men\'eame, Spain's largest social news aggregator, comprising 14,995 articles from 2,868 outlets (2005--2021) with 13.2 million comments from 96,034 users.
Articles are automatically labeled for bias using a supervised classifier trained on the MBBMD corpus, and enriched with 38 interaction features spanning karma distributions, comment engagement, and temporal dynamics.
Our analysis reveals that biased articles exhibit statistically significant differences in user engagement: higher karma entropy ($+4.6\%$, $p < 10^{-24}$), more negative comment sentiment ($59.7\%$ vs.\ $56.2\%$ negative, $p < 10^{-5}$), faster initial reactions ($-10.7\%$ time to 10th comment), and stronger karma degradation over time.
Community detection on the user-outlet bipartite graph identifies two distinct media ecosystems with different bias profiles.
Gradient Boosting achieves the best predictive performance (AUC $0.642$), with article text length and comment karma statistics as the most informative features.
These findings suggest that crowd behavioral signals provide a viable, scalable complement to content-based bias detection, enabling interaction-aware approaches to media bias monitoring.
The dataset and code are publicly available.\footnote{\url{https://doi.org/10.5281/zenodo.19164549}}
\end{abstract}
